\section{Transition to fixed OpenOOD evaluation}
\label{sec:openood-transition}

\headingindent
Near the end of the internship, we replaced the evaluation based on standard
\code{torchvision} test splits with the fixed OpenOOD v1.5 image-list
manifests \cite{yang2022openood,zhang2023openood}. As summarized in
Table~\ref{tab:dataset-protocols}, this transition added the benchmark's
Near-OOD and Far-OOD datasets and fixed the exact images used for every
comparison. For CIFAR-10 and CIFAR-100, the manifests add TIN as Near-OOD and
Textures and Places365 as Far-OOD. For ImageNet-200, they specify SSB-hard and
NINCO as Near-OOD and curated subsets of iNaturalist, Textures, and OpenImage-O
as Far-OOD. This standardized protocol prevents accidental changes in dataset
size and makes results more directly comparable with published OpenOOD values.

While moving to this evaluation, we found that the metric convention was still
not fully standardized. In particular, papers and even different OpenOOD
artifacts alternate between treating ID or OOD samples as the positive class.
The convention does not change AUROC when labels and score direction are
reversed together, but it changes the operating point used by FPR@95. One
convention can therefore produce a substantially lower FPR@95 for exactly the
same predictions without representing a better detector.

\subsection{Two definitions of FPR@95}

\headingindent
Let $s_{\mathrm{ID}}$ increase with ID-likeness and let
$s_{\mathrm{OOD}}=-s_{\mathrm{ID}}$. Two different operating points are
called FPR@95 in the literature.

\begin{table}[H]
\centering
\rowcolors{2}{rowblue!35}{white}
\begin{tabularx}{\textwidth}{@{}L{1.6cm}L{2.5cm}YY@{}}
\toprule
Name & Positive / negative & Threshold condition & Reported false-positive rate \\
\midrule
\idpos & ID / OOD &
$\Pr(s_{\mathrm{ID}}\ge\tau_{\mathrm{ID}}\mid\mathrm{ID})=0.95$ &
$\fprid=\Pr(s_{\mathrm{ID}}\ge\tau_{\mathrm{ID}}\mid\mathrm{OOD})$ \\
\oodpos & OOD / ID &
$\Pr(s_{\mathrm{OOD}}\ge\tau_{\mathrm{OOD}}\mid\mathrm{OOD})=0.95$ &
$\fprood=\Pr(s_{\mathrm{OOD}}\ge\tau_{\mathrm{OOD}}\mid\mathrm{ID})$ \\
\bottomrule
\end{tabularx}
\caption{The two operating points commonly called FPR@95.}
\label{tab:fpr-conventions}
\end{table}

Under \idpos, FPR@95 is the share of OOD inputs accepted as ID while 95\% of
ID inputs are retained. Under \oodpos, it is the share of ID inputs flagged as
OOD while 95\% of OOD inputs are detected. They use different thresholds and
are not complements.

AUROC is not influenced by the convention because it measures ranking. Let
$X_{\mathrm{ID}}$ and $X_{\mathrm{OOD}}$ be independently sampled ID and OOD
examples. Under \idpos, a score $s_{\mathrm{ID}}$ increases with ID-likeness,
\[
  \operatorname{AUROC}_{\mathrm{ID+}}
  = \Pr\!\left(s_{\mathrm{ID}}(X_{\mathrm{ID}})
    > s_{\mathrm{ID}}(X_{\mathrm{OOD}})\right).
\]
Under \oodpos, the score direction is reversed,
$s_{\mathrm{OOD}}=-s_{\mathrm{ID}}$, and the OOD example is the positive one.
Therefore,
\begin{align*}
  \operatorname{AUROC}_{\mathrm{OOD+}}
  &= \Pr\!\left(s_{\mathrm{OOD}}(X_{\mathrm{OOD}})
    > s_{\mathrm{OOD}}(X_{\mathrm{ID}})\right) \\
  &= \Pr\!\left(-s_{\mathrm{ID}}(X_{\mathrm{OOD}})
    > -s_{\mathrm{ID}}(X_{\mathrm{ID}})\right) \\
  &= \Pr\!\left(s_{\mathrm{ID}}(X_{\mathrm{ID}})
    > s_{\mathrm{ID}}(X_{\mathrm{OOD}})\right)
   = \operatorname{AUROC}_{\mathrm{ID+}}.
\end{align*}

% FPR@95 is asymmetric: it reads the ROC curve at 95\% recall for one class or
% the other. The same score vectors can therefore give the same AUROC and very
% different FPR@95 values.

\subsection{OpenOOD version mismatch}

\headingindent
OpenOOD cannot be assigned one convention without naming the artifact and
date. Its v1.0 paper used \idpos. The v1.5 paper and published spreadsheet use
\oodpos, but the tagged \code{v1.5} evaluator still encoded ID as positive.
The evaluator on the later \code{main} branch switched to OOD as positive.
This history explains why our early project values and the OpenOOD spreadsheet
could have matching AUROC but different FPR@95.

\begin{table}[ht]
\centering
\rowcolors{2}{rowblue!35}{white}
\begin{tabularx}{\textwidth}{@{}L{4.5cm}L{2.0cm}Y@{}}
\toprule
Source & Convention & Consequence \\
\midrule
OpenOOD v1.0 paper & \idpos & -- \\
OpenOOD v1.5 paper & \oodpos & -- \\
GitHub tag \code{v1.5} & \idpos & Tagged code disagrees with the v1.5 paper. \\
Later and current \code{main} & \oodpos & Code follows the v1.5 paper. \\
Published v1.5 spreadsheet & \oodpos & Reference numbers match the OOD-positive calculation. \\
\bottomrule
\end{tabularx}
\caption{OpenOOD convention timeline.}
\label{tab:openood-convention-timeline}
\end{table}

We verified the spreadsheet convention with the three released CIFAR-10
ResNet-18 seeds. Mean AUROC was 87.16 under either simultaneous label and score
reversal. Mean \fprid was 61.01, while mean \fprood was 53.12. The spreadsheet
reported $53.08\pm4.86$, which matches the OOD-positive calculation. Thus, the
apparently lower FPR@95 came from evaluating a different operating point, not
from improved predictions.

\subsection{The problem persists in recent work}

\headingindent
An audit of representative and recent papers found both conventions. Most of
the classical methods used \idpos, sometimes only implicitly. Several recent
methods and the later OpenOOD evaluator used \oodpos. Some papers did not state
enough information to decide.

\begin{table}[H]
\centering
\begin{minipage}{9cm}
\centering
\footnotesize
\rowcolors{2}{rowblue!35}{white}
\begin{tabular}{@{}L{5.8cm}L{2.2cm}@{}}
\toprule
Paper or method & Convention \\
\midrule
MSP \cite{hendrycks2017baseline} & \idpos \\
Mahalanobis \cite{lee2018mahalanobis} & \idpos, implicit \\
Deep $k$-NN \cite{sun2022knn} & \idpos \\
PixMix \cite{hendrycks2022pixmix} & Not specified \\
ReAct \cite{sun2021react} & \idpos \\
ASH \cite{djurisic2023ash} & \idpos \\
InterNeg, CVPR 2026 \cite{xu2026interneg} & \oodpos \\
Extended Logit Normalization, CVPR 2026 \cite{ding2026elogitnorm} & \oodpos \\
RankOOD, CVPR 2026 \cite{denipitiyage2026rankood} & \idpos \\
\bottomrule
\end{tabular}
\caption{Examples from the convention audit in different papers. Recent publication date does not
resolve the ambiguity.}
\label{tab:paper-convention-audit}
\end{minipage}
\end{table}

% \subsection{Reporting policy}

% \headingindent
% This project keeps \idpos as its primary operational convention because its
% scores increase with ID-likeness and its main question is OOD false acceptance
% at fixed ID coverage. Future reports should use the explicit label
% \textbf{FPR@95-ID-TPR}. For external comparisons, they should also compute
% \fprood from the same saved per-sample scores.

% Every result should record the positive class, score direction, threshold
% population, false-positive population, benchmark version, checkpoint,
% preprocessing, image lists, and aggregation rule. AUPR should be labeled
% AUPR-IN or AUPR-OUT. Raw scores and labels should remain available so that
% both FPR conventions can be recomputed without rerunning inference.

\subsection{Full results under both conventions}

\headingindent
Tables~\ref{tab:openood-cifar10-full}--\ref{tab:openood-imagenet200-full}
reproduce the full OpenOOD comparisons from the final experiment recap. Every
cell reports AUROC/FPR@95 as a fraction. Blue variant cells denote \idpos and
red variant cells denote \oodpos. Missing entries correspond to datasets that
were not part of the earlier, non-OpenOOD evaluation.

\clearpage
\begin{landscape}
\begin{table}[H]
\centering
\tiny
\setlength{\tabcolsep}{2.25pt}
\renewcommand{\arraystretch}{1.08}
\resizebox{\linewidth}{!}{%
\begin{tabular}{@{}>{\columncolor{conventionred}\bfseries}L{2.45cm}>{\columncolor{rowblue!20}}c>{\columncolor{rowblue!20}}L{1.65cm}
>{\columncolor{neargreen}}c>{\columncolor{neargreen}}c>{\columncolor{nearaverage}\bfseries}c
>{\columncolor{farorange}}c>{\columncolor{farorange}}c>{\columncolor{farorange}}c>{\columncolor{farorange}}c
>{\columncolor{faraverage}\bfseries}c>{\columncolor{macroblue}\bfseries}c@{}}
\toprule
\cellcolor{rowblue!35}Variant & \cellcolor{rowblue!35}Conv. & \cellcolor{rowblue!35}Selected OOD score &
\multicolumn{3}{>{\columncolor{neargreen}}c}{Near-OOD} &
\multicolumn{5}{>{\columncolor{farorange}}c}{Far-OOD} & \cellcolor{macroblue}Balanced \\
\cmidrule(lr){4-6}\cmidrule(lr){7-11}
\cellcolor{rowblue!35} & \cellcolor{rowblue!35} & \cellcolor{rowblue!35} & CIFAR-100 & TIN & Average & MNIST & Places365 & SVHN & Textures & Average & Macro \\
\midrule
\cellcolor{selectedblue!55}Teacher MSP (R18) & \idpos & MSP & .879/.624 & -- & -- & .918/.556 & -- & .904/.598 & -- & -- & .895/.601 \\
\cellcolor{conventionred}Teacher MSP (R18) & \oodpos & MSP & .880/.437 & .891/.378 & .885/.408 & .925/.230 & .894/.373 & .918/.266 & .885/.412 & .906/.321 & .896/.364 \\
\cellcolor{selectedblue!55}Teacher energy (R18) & \idpos & Energy & .879/.506 & -- & -- & .958/.260 & -- & .918/.412 & -- & -- & .908/.421 \\
\cellcolor{conventionred}Teacher energy (R18) & \oodpos & Energy & .879/.541 & .897/.461 & .888/.501 & .963/.159 & .909/.418 & .935/.298 & .879/.549 & .921/.356 & .905/.429 \\
\cellcolor{selectedblue!55}Teacher MSP (R50) & \idpos & MSP & .873/.628 & -- & -- & .897/.622 & -- & .891/.680 & -- & -- & .884/.639 \\
\cellcolor{conventionred}Teacher MSP (R50) & \oodpos & MSP & .873/.472 & .889/.382 & .881/.427 & .920/.237 & .895/.352 & .889/.304 & .876/.417 & .895/.328 & .888/.377 \\
\cellcolor{selectedblue!55}Teacher energy (R50) & \idpos & Energy & .868/.517 & -- & -- & .932/.408 & -- & .911/.517 & -- & -- & .895/.490 \\
\cellcolor{conventionred}Teacher energy (R50) & \oodpos & Energy & .868/.598 & .896/.475 & .882/.536 & .955/.176 & .913/.392 & .908/.325 & .868/.539 & .911/.358 & .897/.447 \\
\midrule
\cellcolor{selectedblue!55}Affine pixels (R50) & \idpos & Student energy & .871/.552 & -- & -- & .906/.528 & -- & .960/.247 & -- & -- & .902/.469 \\
\cellcolor{conventionred}Affine pixels (R50) & \oodpos & Student energy & .871/.565 & .891/.468 & .881/.517 & .931/.255 & .905/.400 & .961/.146 & .847/.619 & .911/.355 & .896/.436 \\
\cellcolor{selectedblue!55}Channel clipping (R18) & \idpos & Student energy & .843/.554 & -- & -- & .946/.288 & -- & .910/.424 & -- & -- & .886/.455 \\
\cellcolor{conventionred}Channel clipping (R18) & \oodpos & Student energy & .844/.639 & .878/.544 & .861/.592 & .937/.302 & .880/.547 & .907/.480 & .893/.508 & .904/.459 & .883/.525 \\
\cellcolor{selectedblue!55}Feature denoising L2 (R18) & \idpos & Absolute improvement & .542/.913 & -- & -- & .597/.918 & -- & .937/.271 & -- & -- & .655/.754 \\
\cellcolor{conventionred}Feature denoising L2 (R18) & \oodpos & Absolute improvement & .541/.953 & .599/.896 & .570/.924 & .753/.652 & .590/.901 & .958/.211 & .647/.993 & .737/.689 & .653/.807 \\
\cellcolor{selectedblue!55}Feature denoising L3 (R18) & \idpos & Absolute improvement & .841/.662 & -- & -- & .846/.835 & -- & .962/.217 & -- & -- & .872/.594 \\
\cellcolor{conventionred}Feature denoising L3 (R18) & \oodpos & Absolute improvement & .840/.548 & .895/.374 & .868/.461 & .880/.309 & .905/.330 & .975/.123 & .930/.346 & .922/.277 & .895/.369 \\
\cellcolor{selectedblue!55}Feature denoising L4 (R18) & \idpos & Absolute improvement & .773/.740 & -- & -- & .979/.115 & -- & .893/.458 & -- & -- & .854/.514 \\
\cellcolor{conventionred}Feature denoising L4 (R18) & \oodpos & Absolute improvement & .772/.721 & .778/.704 & .775/.713 & .980/.100 & .803/.648 & .912/.413 & .751/.793 & .862/.489 & .818/.601 \\
\midrule
Mahalanobis \cite{lee2018mahalanobis} & \oodpos & MDS & .836/.528 & .848/.470 & .842/.499 & .901/.273 & .849/.477 & .912/.260 & .927/.279 & .897/.322 & .870/.411 \\
$k$-NN \cite{sun2022knn} & \oodpos & KNN & .897/.376 & .916/.304 & .906/.340 & .943/.201 & .918/.304 & .927/.226 & .932/.241 & .930/.243 & .918/.291 \\
PixMix \cite{hendrycks2022pixmix} & \oodpos & MSP & .919/.289 & .943/.210 & .931/.250 & .917/.274 & .941/.211 & .979/.101 & .989/.066 & .957/.163 & .944/.206 \\
ReAct \cite{sun2021react} & \oodpos & ReAct & .859/.674 & .883/.597 & .871/.636 & .928/.338 & .903/.442 & .891/.502 & .894/.514 & .904/.449 & .888/.542 \\
ASH \cite{djurisic2023ash} & \oodpos & ASH & .741/.873 & .764/.863 & .753/.868 & .832/.700 & .799/.779 & .735/.836 & .775/.846 & .785/.790 & .769/.829 \\
ViM \cite{wang2022vim} & \oodpos & ViM & .877/.492 & .896/.405 & .887/.448 & .948/.184 & .895/.414 & .945/.193 & .952/.211 & .935/.251 & .911/.349 \\
\bottomrule
\end{tabular}%
}
\caption{CIFAR-10 as the ID dataset. Cells report AUROC/FPR@95 under the stated positive-class convention.}
\label{tab:openood-cifar10-full}
\end{table}
\end{landscape}

\clearpage
\begin{landscape}
\begin{table}[H]
\centering
\tiny
\setlength{\tabcolsep}{2.25pt}
\renewcommand{\arraystretch}{1.08}
\resizebox{\linewidth}{!}{%
\begin{tabular}{@{}>{\columncolor{conventionred}\bfseries}L{2.45cm}>{\columncolor{rowblue!20}}c>{\columncolor{rowblue!20}}L{1.65cm}
>{\columncolor{neargreen}}c>{\columncolor{neargreen}}c>{\columncolor{nearaverage}\bfseries}c
>{\columncolor{farorange}}c>{\columncolor{farorange}}c>{\columncolor{farorange}}c>{\columncolor{farorange}}c
>{\columncolor{faraverage}\bfseries}c>{\columncolor{macroblue}\bfseries}c@{}}
\toprule
\cellcolor{rowblue!35}Variant & \cellcolor{rowblue!35}Conv. & \cellcolor{rowblue!35}Selected OOD score &
\multicolumn{3}{>{\columncolor{neargreen}}c}{Near-OOD} &
\multicolumn{5}{>{\columncolor{farorange}}c}{Far-OOD} & \cellcolor{macroblue}Balanced \\
\cmidrule(lr){4-6}\cmidrule(lr){7-11}
\cellcolor{rowblue!35} & \cellcolor{rowblue!35} & \cellcolor{rowblue!35} & CIFAR-10 & TIN & Average & MNIST & Places365 & SVHN & Textures & Average & Macro \\
\midrule
\cellcolor{selectedblue!55}Teacher MSP (R18) & \idpos & MSP & .792/.794 & -- & -- & .701/.956 & -- & .806/.812 & -- & -- & .773/.839 \\
\cellcolor{conventionred}Teacher MSP (R18) & \oodpos & MSP & .790/.593 & .833/.484 & .811/.538 & .776/.556 & .802/.545 & .795/.545 & .781/.603 & .788/.562 & .800/.550 \\
\cellcolor{selectedblue!55}Teacher energy (R18) & \idpos & Energy & .795/.800 & -- & -- & .695/.974 & -- & .829/.771 & -- & -- & .778/.836 \\
\cellcolor{conventionred}Teacher energy (R18) & \oodpos & Energy & .793/.597 & .835/.483 & .814/.540 & .780/.553 & .801/.554 & .816/.521 & .789/.605 & .796/.558 & .805/.549 \\
\cellcolor{selectedblue!55}Teacher MSP (R50) & \idpos & MSP & .792/.781 & -- & -- & .811/.823 & -- & .758/.793 & -- & -- & .788/.795 \\
\cellcolor{conventionred}Teacher MSP (R50) & \oodpos & MSP & .790/.618 & .819/.541 & .804/.580 & .820/.437 & .786/.632 & .741/.729 & .766/.675 & .778/.618 & .791/.599 \\
\cellcolor{selectedblue!55}Teacher energy (R50) & \idpos & Energy & .801/.774 & -- & -- & .878/.670 & -- & .783/.766 & -- & -- & .816/.746 \\
\cellcolor{conventionred}Teacher energy (R50) & \oodpos & Energy & .799/.645 & .828/.580 & .813/.613 & .878/.366 & .781/.705 & .773/.712 & .779/.686 & .803/.617 & .808/.615 \\
\midrule
\cellcolor{selectedblue!55}Affine pixels (R50) & \idpos & Student energy & .779/.797 & -- & -- & .875/.698 & -- & .925/.431 & -- & -- & .839/.681 \\
\cellcolor{conventionred}Affine pixels (R50) & \oodpos & Student energy & .777/.715 & .817/.596 & .797/.656 & .886/.334 & .770/.716 & .929/.248 & .797/.623 & .846/.480 & .821/.568 \\
\cellcolor{selectedblue!55}Channel clipping (R18) & \idpos & Student energy & .751/.842 & -- & -- & .797/.879 & -- & .940/.340 & -- & -- & .810/.726 \\
\cellcolor{conventionred}Channel clipping (R18) & \oodpos & Student energy & .749/.699 & .829/.508 & .789/.604 & .865/.392 & .807/.575 & .937/.256 & .858/.547 & .867/.442 & .828/.523 \\
\cellcolor{selectedblue!55}Feature denoising L2 (R18) & \idpos & Absolute improvement & .552/.950 & -- & -- & .881/.545 & -- & .900/.421 & -- & -- & .721/.717 \\
\cellcolor{conventionred}Feature denoising L2 (R18) & \oodpos & Absolute improvement & .553/.894 & .536/.898 & .545/.896 & .916/.322 & .571/.879 & .914/.372 & .408/.999 & .702/.643 & .623/.770 \\
\cellcolor{selectedblue!55}Feature denoising L3 (R18) & \idpos & Absolute improvement & .504/.967 & -- & -- & .994/.022 & -- & .936/.283 & -- & -- & .734/.560 \\
\cellcolor{conventionred}Feature denoising L3 (R18) & \oodpos & Absolute improvement & .505/.930 & .556/.892 & .531/.911 & .988/.053 & .616/.836 & .928/.353 & .613/.989 & .786/.558 & .658/.734 \\
\cellcolor{selectedblue!55}Feature denoising L4 (R18) & \idpos & Absolute improvement & .774/.825 & -- & -- & .621/.981 & -- & .806/.760 & -- & -- & .744/.848 \\
\cellcolor{conventionred}Feature denoising L4 (R18) & \oodpos & Absolute improvement & .769/.631 & .780/.568 & .775/.599 & .716/.646 & .758/.631 & .784/.579 & .743/.696 & .750/.638 & .762/.619 \\
\midrule
Mahalanobis \cite{lee2018mahalanobis} & \oodpos & MDS & .559/.880 & .615/.790 & .587/.835 & .675/.717 & .631/.796 & .707/.672 & .763/.705 & .694/.723 & .640/.779 \\
$k$-NN \cite{sun2022knn} & \oodpos & KNN & .770/.728 & .833/.496 & .802/.612 & .824/.486 & .794/.607 & .842/.517 & .837/.536 & .824/.536 & .813/.574 \\
PixMix \cite{hendrycks2022pixmix} & \oodpos & MSP & .766/.622 & .832/.515 & .799/.568 & .696/.703 & .814/.551 & .934/.308 & .918/.375 & .841/.484 & .820/.526 \\
ReAct \cite{sun2021react} & \oodpos & ReAct & .787/.613 & .829/.515 & .808/.564 & .784/.560 & .800/.553 & .830/.504 & .802/.550 & .804/.542 & .806/.553 \\
ASH \cite{djurisic2023ash} & \oodpos & ASH & .765/.681 & .799/.634 & .782/.657 & .772/.666 & .788/.630 & .856/.460 & .807/.613 & .806/.592 & .794/.625 \\
ViM \cite{wang2022vim} & \oodpos & ViM & .722/.706 & .778/.547 & .750/.626 & .819/.483 & .758/.616 & .831/.462 & .859/.469 & .817/.507 & .783/.567 \\
\bottomrule
\end{tabular}%
}
\caption{CIFAR-100 as the ID dataset. Cells report AUROC/FPR@95 under the stated positive-class convention.}
\label{tab:openood-cifar100-full}
\end{table}
\end{landscape}

\clearpage
\begin{landscape}
\begin{table}[H]
\centering
\tiny
\setlength{\tabcolsep}{4pt}
\renewcommand{\arraystretch}{1.08}
\resizebox{\linewidth}{!}{%
\begin{tabular}{@{}>{\columncolor{conventionred}\bfseries}L{2.5cm}>{\columncolor{rowblue!20}}c>{\columncolor{rowblue!20}}L{1.75cm}
>{\columncolor{neargreen}}c>{\columncolor{neargreen}}c>{\columncolor{nearaverage}\bfseries}c
>{\columncolor{farorange}}c>{\columncolor{farorange}}c>{\columncolor{farorange}}c
>{\columncolor{faraverage}\bfseries}c>{\columncolor{macroblue}\bfseries}c@{}}
\toprule
\cellcolor{rowblue!35}Variant & \cellcolor{rowblue!35}Conv. & \cellcolor{rowblue!35}Selected OOD score &
\multicolumn{3}{>{\columncolor{neargreen}}c}{Near-OOD} &
\multicolumn{4}{>{\columncolor{farorange}}c}{Far-OOD} & \cellcolor{macroblue}Balanced \\
\cmidrule(lr){4-6}\cmidrule(lr){7-10}
\cellcolor{rowblue!35} & \cellcolor{rowblue!35} & \cellcolor{rowblue!35} & SSB-hard & NINCO & Average & iNaturalist & Textures & OpenImage-O & Average & Macro \\
\midrule
\cellcolor{selectedblue!55}Teacher MSP & \idpos & MSP & .804/.726 & .862/.648 & .833/.687 & .931/.398 & .883/.496 & .892/.564 & .902/.486 & .867/.586 \\
\cellcolor{conventionred}Teacher MSP & \oodpos & MSP & .804/.660 & .862/.445 & .833/.552 & .931/.257 & .883/.453 & .892/.353 & .902/.355 & .867/.454 \\
\cellcolor{selectedblue!55}Teacher energy & \idpos & Energy & .798/.730 & .853/.661 & .826/.695 & .932/.414 & .906/.413 & .896/.558 & .911/.462 & .868/.579 \\
\cellcolor{conventionred}Teacher energy & \oodpos & Energy & .798/.695 & .853/.494 & .826/.595 & .932/.237 & .906/.429 & .896/.355 & .911/.340 & .868/.467 \\
\midrule
\cellcolor{selectedblue!55}Affine pixels & \idpos & Student energy & .791/.736 & .844/.685 & .817/.711 & .921/.473 & .907/.410 & .889/.578 & .906/.487 & .862/.599 \\
\cellcolor{conventionred}Affine pixels & \oodpos & Student energy & .791/.712 & .844/.512 & .817/.612 & .921/.270 & .907/.414 & .889/.382 & .906/.355 & .862/.484 \\
\cellcolor{selectedblue!55}Channel clipping & \idpos & Student energy & .770/.760 & .830/.695 & .800/.728 & .919/.450 & .913/.389 & .877/.605 & .903/.481 & .852/.604 \\
\cellcolor{conventionred}Channel clipping & \oodpos & Student energy & .770/.735 & .830/.568 & .800/.651 & .919/.291 & .913/.372 & .877/.401 & .903/.355 & .852/.503 \\
\cellcolor{selectedblue!55}Spatial clipping & \idpos & Energy gap & .797/.734 & .852/.664 & .825/.699 & .931/.423 & .905/.420 & .895/.563 & .910/.469 & .868/.584 \\
\cellcolor{conventionred}Spatial clipping & \oodpos & Energy gap & .797/.695 & .852/.495 & .825/.595 & .931/.238 & .905/.429 & .895/.355 & .910/.341 & .868/.468 \\
\cellcolor{selectedblue!55}Feature denoising, layer 2 & \idpos & Absolute improvement & .628/.896 & .617/.843 & .623/.869 & .446/.984 & .608/.782 & .489/.915 & .515/.894 & .569/.882 \\
\cellcolor{conventionred}Feature denoising, layer 2 & \oodpos & Absolute improvement & .628/.886 & .617/.923 & .623/.904 & .446/.938 & .608/.997 & .489/.973 & .515/.970 & .569/.937 \\
\cellcolor{selectedblue!55}Feature denoising, layer 3 & \idpos & Absolute improvement & .642/.894 & .695/.797 & .668/.845 & .745/.733 & .707/.673 & .621/.849 & .691/.752 & .680/.799 \\
\cellcolor{conventionred}Feature denoising, layer 3 & \oodpos & Absolute improvement & .642/.848 & .695/.826 & .668/.837 & .745/.766 & .707/.978 & .621/.909 & .691/.885 & .680/.861 \\
\cellcolor{selectedblue!55}Feature denoising, layer 4 & \idpos & Absolute improvement & .722/.797 & .761/.755 & .742/.776 & .869/.556 & .908/.367 & .847/.607 & .875/.510 & .808/.643 \\
\cellcolor{conventionred}Feature denoising, layer 4 & \oodpos & Absolute improvement & .723/.782 & .762/.713 & .742/.747 & .871/.474 & .908/.444 & .847/.540 & .875/.486 & .809/.617 \\
\midrule
Mahalanobis \cite{lee2018mahalanobis} & \oodpos & MDS & .584/.837 & .655/.746 & .619/.791 & .750/.585 & .792/.582 & .699/.683 & .747/.617 & .683/.704 \\
$k$-NN \cite{sun2022knn} & \oodpos & KNN & .770/.737 & .861/.466 & .816/.602 & .940/.245 & .953/.244 & .902/.329 & .932/.273 & .874/.437 \\
PixMix \cite{hendrycks2022pixmix} & \oodpos & MSP & .788/.678 & .855/.463 & .822/.570 & .925/.273 & .898/.444 & .884/.371 & .902/.363 & .862/.467 \\
ReAct \cite{sun2021react} & \oodpos & ReAct & .790/.715 & .848/.535 & .819/.625 & .937/.230 & .929/.297 & .904/.329 & .923/.285 & .871/.455 \\
ASH \cite{djurisic2023ash} & \oodpos & ASH & .795/.721 & .852/.576 & .824/.649 & .951/.225 & .948/.257 & .918/.337 & .939/.273 & .881/.461 \\
ViM \cite{wang2022vim} & \oodpos & ViM & .740/.713 & .833/.471 & .787/.592 & .910/.273 & .946/.204 & .882/.339 & .913/.272 & .850/.432 \\
\bottomrule
\end{tabular}%
}
\caption{Selected ResNet-18 variants and published OpenOOD baselines on ImageNet-200. Cells report AUROC/FPR@95 under the stated positive-class convention.}
\label{tab:openood-imagenet200-full}
\end{table}
\end{landscape}

